El estudio de exoplanetas plantea un problema de estructura: no siempre esperamos encontrar grupos simples, separados y perfectamente definidos. Las propiedades planetarias cambian de forma continua. Un planeta puede ser pequeño pero muy irradiado, masivo pero lejano, denso pero con una órbita particular, o pertenecer a una región del catálogo que está sobrerrepresentada por un método de descubrimiento específico. Por eso, antes de asumir clases planetarias rígidas, resulta útil preguntarse si el catálogo contiene regiones conectadas, transiciones, ramas o ciclos que puedan ser inspeccionados desde una perspectiva astrofísica.

En este proyecto usamos Mapper, una técnica de Topological Data Analysis, para construir grafos exploratorios a partir de distintas selecciones de variables de exoplanetas. A diferencia de un clustering tradicional, Mapper no fuerza una partición única del catálogo. En lugar de asignar cada planeta a una sola clase, construye un grafo de regiones locales que pueden solaparse y conectarse entre sí. Esta propiedad es importante para este problema, porque muchas poblaciones planetarias no necesariamente aparecen como grupos compactos, sino como continuos o transiciones entre regímenes físicos, orbitales y térmicos.

La pregunta de investigación que guía el reporte es la siguiente:

\begin{quote}
¿Pueden los grafos Mapper construidos sobre variables físicas, orbitales, térmicas y combinadas de exoplanetas revelar regiones candidatas para inspección científica, y en qué medida esas regiones pueden confundirse con sesgo observacional, dependencia de imputación o sensibilidad metodológica?
\end{quote}

La lectura del proyecto es deliberadamente cautelosa. Mapper no se interpreta aquí como una prueba definitiva de estructura física ni como una taxonomía final de exoplanetas. En este trabajo, una ``estructura'' significa una organización medible dentro del grafo: número de nodos, aristas, componentes, ciclos, regiones con coherencia física u orbital, regiones dominadas por un método de descubrimiento, o regiones con alta o baja dependencia de imputación. En otras palabras, el grafo se entiende como un resumen exploratorio de vecindades inducidas por una matriz de datos procesada y completada, no como la topología real de la población planetaria.

El análisis compara varias maneras de mirar el mismo catálogo. El espacio físico aproxima qué tipo de planeta es un objeto mediante variables como radio, masa y densidad. El espacio orbital describe dónde vive y cómo se organiza su órbita mediante periodo orbital y semieje mayor. El espacio térmico aproxima el ambiente energético mediante insolación y temperatura de equilibrio. Finalmente, los espacios conjuntos permiten observar cómo se combinan estas descripciones. Esta separación es central porque evita mezclar todas las variables desde el inicio y permite preguntar qué tipo de estructura produce cada perspectiva.

A partir de esta pregunta general, el reporte evalúa cuatro hipótesis de trabajo. Primero, una hipótesis física u orbital: algunos espacios de variables, especialmente el orbital, podrían producir grafos con estructura topológica rica y regiones interpretables. Segundo, una hipótesis de imputación: las estructuras que dependen fuertemente de variables imputadas deben considerarse menos confiables. Tercero, una hipótesis de sesgo observacional: algunas regiones Mapper podrían estar enriquecidas por \texttt{discoverymethod}, \texttt{disc\_year} o \texttt{disc\_facility}, lo que indicaría que parte de la topología refleja el proceso de descubrimiento más que una diferencia física intrínseca. Cuarto, una hipótesis mixta: la estructura observada probablemente no será puramente física ni puramente observacional, sino una combinación de señal astrofísica, sesgo de detección y decisiones metodológicas.

Esta distinción es importante porque una región del grafo puede parecer científicamente interesante por varias razones distintas. Puede ser compatible con una señal astrofísica si muestra coherencia en radio, masa, densidad, periodo orbital, insolación o temperatura. Puede ser metodológicamente confiable si aparece con baja imputación y no depende de una sola construcción del lente. Pero también puede ser observacionalmente sospechosa si está dominada por un solo método de descubrimiento, por una facilidad específica o por una época de observación. Por eso, en lugar de concluir directamente que una rama o componente representa una familia planetaria, el análisis pregunta qué tipo de evidencia sostiene esa región.

El objetivo del reporte, entonces, no es decir simplemente si ``hay clases'' o ``no hay clases''. El objetivo es descomponer la topología observada en varias fuentes de evidencia:

\[
\text{topología observada}
=
\text{señal física/orbital}
+
\text{sesgo observacional}
+
\text{efecto de imputación}
+
\text{sensibilidad metodológica}.
\]

Bajo esta formulación, una conclusión fuerte no sería afirmar que Mapper descubre una clasificación definitiva de exoplanetas, sino identificar qué regiones son candidatas a interpretación física, cuáles parecen dominadas por sesgo observacional, cuáles mezclan ambos efectos y cuáles deben considerarse débiles. Esta lectura permite aprovechar la riqueza del catálogo sin ignorar las limitaciones de los datos con los que fue construido.